\section{Engineering observations}
\label{sec:engineering}

\subsection{Accelerated runtime}

The run used an NVIDIA A100 80GB PCIe on RunPod Secure Cloud with CUDA 13.0 and
BF16 support. The locked environment recorded Python 3.12.13, PyTorch 2.13.0,
Transformers 5.14.1, TRL 1.9.2, PEFT 0.20.0, Flash Linear Attention 0.5.2, and
\Code{causal-conv1d} 1.7.0. Cached preparation of the latter took roughly 1.1
seconds. The much longer first preflight delay came from Triton compilation and
autotuning of the gated-delta path; FLA documents its persistent autotune cache
controls \citep{flashlinearattention}.

Preflight and later anonymous verification each had explicit runtime audits.
The final receipt records a \emph{two-token non-generative forward probe} that
observed one call to \Code{causal\_conv1d\_fn} and one call to
\Code{chunk\_gated\_delta\_rule}, followed by CUDA synchronization.
\claimsource{publication} This establishes that both required accelerated paths
could execute in the verified environment. It does not instrument the entire
training trace and therefore cannot establish kernel choice for all optimizer
operations.

\subsection{Host layout, persistence, and retrieval}

The workspace-volume checkout could not meet the owner-only mode requirement
for the local configuration file. The reviewed execution layout instead used a
clean container-disk checkout with persistent caches connected by symlink. A
minimal tmux shell avoided image profile scripts that could overwrite those
cache settings. Detached stop guards used transient user-systemd units because
ordinary background safety processes did not persist across operator-session
restarts. These corrections were completed before the untouched-base run at
source commit \Code{8645addf427edf7ac218ed977a0be9102342851f}.
\claimsource{metadata}

Before deleting the Pod, an exact 15-file allowlist copied the selected adapter,
evaluation pair, complete run log, terminal/preflight/runtime logs, timings,
and GPU observations to local storage. A supplemental archive retained both
\Code{checkpoint-84} and \Code{checkpoint-210}. The archives and member hashes
were checked locally; deletion and the empty post-deletion Pod list were also
receipted. The final evidence states that no GPU dependency remains.
\claimsource{metadata}

\subsection{Public adapter verification}

The selected adapter was published in the public repository linked in
Section~\ref{sec:reproducibility} at immutable revision
\Fingerprint{dd0ded7bbb5231f204deff9acc63089f4bb5178d}
An anonymous A100 load against the exact base revision returned
``rainbow unicorn.'' for the requested description. This was a
\emph{single-prompt smoke test}, not a replay of the complete final suite.
Local finalization checked the publication receipt and added exactly that model
to the dedicated Collection linked in Section~\ref{sec:reproducibility}.
\claimsource{publication}
